\documentclass[12pt,a4paper,oneside]{report}
\usepackage[T1]{fontenc}
\usepackage[utf8]{inputenc}
\usepackage[french]{babel}
\usepackage{lmodern}
\usepackage{microtype}
\usepackage{setspace}
\onehalfspacing
\setlength{\parindent}{0pt}
\setlength{\parskip}{0.5\baselineskip}
\setlength{\emergencystretch}{3em}
% Prose paragraphs in this manual are short enough to fit on a page.
% Move the whole paragraph rather than splitting instructions across pages.
\interlinepenalty=10000
\clubpenalty=10000
\widowpenalty=10000
\displaywidowpenalty=10000
\raggedbottom
\usepackage[a4paper,margin=2.4cm]{geometry}
\usepackage{array,longtable,booktabs}
\usepackage{etoolbox}
\newcolumntype{L}[1]{>{\raggedright\arraybackslash}p{#1}}
\AtBeginEnvironment{longtable}{\setlength{\parskip}{0pt}}
\usepackage{xcolor}
\usepackage{tikz}
\usepackage{hyperref}
\hypersetup{colorlinks=true,linkcolor=blue!45!black,urlcolor=blue!45!black,
 pdftitle={IPR -- Livre utilisateur},pdfauthor={Projet IPR},pageanchor=false}
\newcommand{\touche}[1]{\textsf{#1}}
\newcommand{\chemin}[1]{\texttt{\detokenize{#1}}}
\newcommand{\attention}[1]{\par\noindent\textbf{À retenir.} #1\par}
\newcommand{\gestesreglages}{\paragraph{Gestes et clavier.}
Cliquer sur le curseur pour lui donner le focus, puis utiliser les flèches
pour avancer au pas prévu par ce réglage. Le champ numérique permet une saisie
directe à confirmer avec \touche{Entrée}. Double-cliquer sur le nom ou le curseur,
ou cliquer sur sa flèche de réinitialisation, restaure sa référence.
Le clic droit sur le nom ou le curseur ouvre aussi la remise à zéro et la saisie.
Ce double-clic n’a pas la même fonction dans le graphe des courbes.}
\begin{document}
\begin{titlepage}
\centering
\vspace*{2cm}
{\Huge\bfseries IPR\par}
\vspace{0.5cm}
{\LARGE Livre utilisateur\par}
\vspace{1cm}
{\large Ouvrir, développer, comparer et exporter ses photos\par}
\vfill
Édition française -- prototype alpha\par
Interface intégrée dans \texttt{dev} (\texttt{a15f703})\par
État relu : 22 octobre 2026\par
\vfill
\begin{minipage}{0.88\textwidth}
Ce livre décrit les commandes présentes, leurs gestes d’utilisation et leurs
limites. Il ne constitue ni une certification de rendu pour tous les boîtiers,
ni la description de fonctions futures. Les chapitres réservés sont volontairement vides.
\end{minipage}
\end{titlepage}
\pagenumbering{roman}
\hypersetup{pageanchor=true}
\tableofcontents
\clearpage
\pagenumbering{arabic}
\part{Prendre en main IPR}
\chapter{Première ouverture et repères}
\label{chap:demarrage}
\section{Ouvrir une photo}
Depuis l’écran d’accueil, cliquer sur \emph{Ouvrir…} et sélectionner une image.
IPR ouvre son dossier contenant et affiche ses images dans le ruban du bas.
Les menus \emph{Fichier / Ouvrir une image…} et \emph{Ouvrir un dossier…}
permettent ensuite de changer de photo ou de dossier. Aucun raccourci dédié
à l’ouverture n’est actuellement déclaré dans l’application.

Un RAW contient des données à développer ; un JPEG, PNG ou TIFF contient déjà
une image rendue. Un DNG peut contenir des données de natures différentes :
son extension ne signifie pas à elle seule qu’il s’agit d’un RAW mosaïqué.
Utiliser les filtres du dialogue pour connaître les fichiers proposés par le build installé.

\section{Les quatre zones de travail}
\begin{center}
\begin{tikzpicture}[x=1cm,y=1cm,every node/.style={font=\small,align=center}]
\draw[fill=gray!8] (0,1) rectangle (2.6,5);
\draw[fill=gray!3] (2.8,1) rectangle (10,5);
\draw[fill=gray!8] (10.2,1) rectangle (13.2,5);
\draw[fill=gray!8] (0,0) rectangle (13.2,0.8);
\node at (1.3,3) {Informations\\Métadonnées\\Notation};
\node at (6.4,3) {Zone image\\Zoom et comparaison};
\node at (11.7,3) {Réglages\\Histogramme fixe\\Familles\\repliables};
\node at (6.6,0.4) {Ruban des images du dossier};
\end{tikzpicture}
\end{center}
Ce schéma représente l’organisation de l’interface, pas une capture d’écran.
Le menu \emph{Affichage} permet de masquer les informations, l’histogramme
ou le ruban. Un clic sur l’en-tête d’une famille déplie ou replie ses réglages.
Le repliement ne remet pas les valeurs à zéro.

\section{Un premier développement simple}
\begin{enumerate}
\item Ouvrir une photo et attendre la fin du traitement.
\item Vérifier son profil, puis ajuster l’exposition dans \emph{Lumière}.
\item Comparer l’original et l’état courant ; vérifier les détails au zoom.
\item Ajuster la couleur, puis le bruit et la netteté si nécessaire.
\item Exporter une copie avec le récapitulatif décrit au chapitre~\ref{chap:export}.
\end{enumerate}

\section{Comprendre les traitements en cours}
Un badge apparaît sur la zone image lorsqu’un traitement dure suffisamment
longtemps. Un aperçu boîtier peut précéder le développement IPR : il est alors
identifié explicitement. La description RAW / Profil reste stable pendant
un simple calcul de retouche pour ne pas déplacer les commandes.
Le message au bas de la fenêtre apporte les détails et les erreurs.

\attention{Le badge est générique, pas une jauge en pourcentage. Il n’expose
pas encore séparément toutes les étapes du moteur. L’absence de badge ne
remplace pas la lecture d’un éventuel message d’erreur.}

\section{Retouches et sauvegarde}
Les réglages sont mémorisés dans un fichier compagnon \chemin{.ipr}, à côté
de la source : par exemple \chemin{photo.CR3.ipr}. Le fichier source n’est
pas réécrit par ces retouches. Garder le compagnon avec l’image lors d’un
déplacement ou d’une sauvegarde ; vérifier les droits d’écriture du dossier.
Le stockage est automatique et différé : attendre les traitements et vérifier
les éventuels avertissements de sauvegarde avant de quitter.

\attention{Le fichier compagnon n’est ni une copie de la photo ni un export.
Les préférences de l’interface sont, elles, stockées dans un INI distinct.}

\chapter{Navigation, zoom et comparaison}
\label{chap:navigation}
\section{Passer d’une image à l’autre}
Cliquer sur une vignette pour ouvrir sa photo. Le contour de la vignette et
son nom en gras identifient l’image effectivement ouverte, indépendamment des
cases de sélection multiple. Si le ruban est replié, les boutons précédent
et suivant apparaissent en bas à droite de la zone image. Leur survol indique
le nom de la cible ; ils ne dépassent pas les extrémités du dossier.

Le bouton centré au-dessus du ruban le replie. Lorsqu’il est masqué, ce bouton
revient en bas de la vue et permet de le rouvrir. Son état est mémorisé.
Les flèches du clavier ne sont pas des raccourcis globaux de changement de photo.

\section{Examiner les détails}
Cliquer dans la photo pour passer du mode ajusté à un agrandissement, puis
cliquer de nouveau sans déplacer la souris pour revenir au mode ajusté.
Quand l’image dépasse la vue, maintenir le bouton gauche et déplacer la souris
pour parcourir les détails. Un clic dans une marge vide ne déclenche pas le zoom.

\paragraph{Clavier et molette.}
\touche{Ctrl + +} agrandit ; \touche{Ctrl + $-$} réduit.
\touche{Ctrl + 0} ajuste à la vue et \touche{Ctrl + 1} affiche à 100\,\%.
\touche{Ctrl + molette} zoome autour de la souris. Sans Ctrl, la molette déplace
l’image agrandie ; \touche{Maj + molette verticale} permet un déplacement horizontal.
Les boutons correspondants restent disponibles au-dessus de la photo.

\attention{100\,\% signifie un pixel de l’aperçu par unité de la vue ;
l’aperçu interactif peut être réduit par rapport à la source pleine résolution.
L’export est reconstruit indépendamment à sa résolution prévue.}

\section{Afficher un comparatif partagé}
Cliquer sur \emph{Avant/après} ou utiliser \touche{Ctrl + Maj + C} lorsque
les deux aperçus sont prêts. L’origine est à gauche et l’état courant à droite.
Déplacer le séparateur pour choisir la largeur des deux zones. Les deux côtés
partagent le même zoom et le même déplacement.

\section{Voir temporairement l’original}
Donner le focus à la vue image, puis maintenir \touche{Espace}. Relâcher revient
à la présentation précédente, y compris au comparatif partagé s’il était actif.
Le bouton \emph{Origine} remplit le même rôle tant que le clic est maintenu.
La commande ne s’applique pas à une saisie dans un champ numérique.
Un changement de photo réinitialise la vue et la comparaison.

\attention{L’origine est la référence de développement exposée par IPR,
pas la mosaïque RAW brute ni nécessairement le JPEG du boîtier. Modifier
le profil de développement peut modifier cette référence.}

\section{Menu contextuel de la photo}
Un clic droit dans la vue ouvre les actions de la photo : export et
réinitialisation des retouches, ainsi que les commandes de fichier.
Le clic droit ne déclenche pas le zoom. Avec le focus dans la vue,
\touche{Maj + F10} ou la touche Menu ouvre aussi ce contexte.
La réinitialisation demande confirmation ; elle n’a pas d’annulation générale.
Les actions Enregistrer et Dupliquer restent désactivées tant que leur
raccordement moteur manque.

\chapter{Informations, classement et sélection}
\section{Lire les informations}
Le panneau gauche décrit la photo ouverte : nom, date disponible, stockage,
dimensions, appareil et autres métadonnées connues. Une information absente
ne doit pas être interprétée comme une valeur mesurée égale à zéro.
Sous l’histogramme, les informations de prise de vue disponibles facilitent
la lecture des ISO, de la focale et du temps d’exposition.

\section{Noter et marquer une photo}
Cliquer sur une des cinq étoiles pour donner une note de 1 à 5.
Cliquer sur la note déjà sélectionnée la remet à zéro. Le drapeau se bascule
par un clic. Choisir un libellé rouge, jaune, vert, bleu ou violet ; cliquer
sur la même couleur retire le libellé. Ces choix de classement ne recolorent
pas les pixels et ne correspondent pas aux bandes du mélangeur.

Aucun raccourci numérique de notation ou de drapeau n’est déclaré pour
l’instant. Les commandes se trouvent dans le panneau Informations.

\section{Copier les retouches à plusieurs photos}
\begin{enumerate}
\item Ouvrir la photo portant les réglages à copier.
\item Cocher les cases des vignettes destinataires, sans confondre case cochée
et image ouverte. Le compteur indique le nombre de fichiers sélectionnés.
\item Cliquer sur \emph{Appliquer à N}. Vérifier le message de résultat.
\item Décocher les fichiers, ou utiliser le bouton de désélection générale.
\end{enumerate}
Cette action écrit les retouches dans les compagnons des destinataires.
Elle n’est pas limitée à la seule famille actuellement dépliée. Les profils
RAW et la balance automatique propres aux destinataires ne sont pas copiés
aveuglément entre boîtiers ; leurs métadonnées de classement sont conservées.

\attention{Il s’agit d’une application groupée des retouches, pas d’une commande
d’export multiple. Il n’existe pas de recette utilisateur de sélection par
Ctrl-clic ou Maj-clic dans le ruban actuel : utiliser les cases.}

\section{Les badges du ruban}
Le badge de la vignette ouverte peut rappeler sa note, son drapeau et son
libellé. L’absence de badge sur les autres vignettes ne signifie pas qu’elles
ne sont pas notées ou retouchées : ces données ne sont pas encore exposées
pour toutes les lignes du modèle.

\section{Clic droit sur une vignette}
Le menu vise la vignette sous la souris. Sur une photo différente de celle
ouverte, choisir d’abord \emph{Ouvrir cette image} : le clic droit seul ne
la charge pas et ne change pas la sélection multiple. Les actions de retouche
restent désactivées pour cette cible tant qu’elle n’est pas ouverte, afin
de ne pas réinitialiser ou exporter la photo précédente par erreur.

\emph{Dupliquer le fichier et ses réglages}, associé au raccourci réservé
\touche{Ctrl + D}, doit créer une copie physique et son compagnon cohérent,
pas un nouveau TIFF développé ni une copie virtuelle. Cette commande attend
encore l’interface moteur et reste désactivée dans la version décrite.

\part{Familles de réglages déjà intégrées}
\chapter{Développement RAW et profil}
\section{Choisir le rendu de départ}
Le groupe \emph{Développement RAW / Profil} permet de choisir un profil DCP
compatible avec la photo. Les propositions dépendent du boîtier, des profils
disponibles et du build. Vérifier la description du profil effectivement utilisé,
pas seulement le nom demandé dans une liste.

Un profil appareil définit un rendu de départ : il peut changer les couleurs,
le contraste et la luminosité apparente avant les retouches. Deux logiciels
ne produisent pas nécessairement les mêmes pixels avec un intitulé similaire.
Ce groupe n’applique pas un DCP d’appareil à un JPEG comme s’il s’agissait d’un RAW.

\section{Changer de profil}
Choisir une entrée dans la liste, attendre le calcul puis comparer le rendu.
Le moteur peut éviter un redécodage complet lorsqu’un tampon compatible est
disponible ; l’interface indique cette capacité. Elle ne garantit pas que
tous les changements de développement sont instantanés.
Les commandes susceptibles de modifier le développement sont désactivées
pendant certains traitements.

\section{Balance automatique et calibration}
La case \emph{Balance des blancs automatique} se trouve ici, une seule fois.
Les réglages manuels de température et de nuance restent dans \emph{Couleur}.
Ne pas confondre balance automatique, température de prise de vue et profil DCP.

La calibration RAW est une option avancée des préférences, accessibles avec
\touche{Ctrl + virgule}, onglet \emph{Avancé · RAW}. Sur les boîtiers et modes
compatibles, le moteur dispose d’une calibration par défaut. Vérifier son état
appliqué dans la description : une case activée n’assure pas qu’un fichier
inconnu possède une caractérisation disponible.

\attention{Le choix avancé concerne l’image ouverte et son compagnon,
pas une nouvelle préférence globale pour toutes les prochaines photos.
Ne changer le manifeste de calibration que dans un protocole maîtrisé.}

\paragraph{Exemple.}
Pour comparer deux profils, garder les retouches identiques, désactiver les
corrections non nécessaires au test, changer uniquement le profil et examiner
les tons de peau, les couleurs vives et les hautes lumières. Noter la balance
des blancs choisie pour rendre la comparaison reproductible.

\chapter{Lumière et histogramme}
\section{Lire l’histogramme fixe}
L’histogramme reste en haut du panneau Réglages pendant le défilement.
Il décrit l’aperçu sRGB courant après traitement, pas les valeurs brutes du
capteur. Les indicateurs d’extrémité signalent des valeurs de l’aperçu à ses
bornes ; ils ne mesurent pas directement la saturation du capteur.

Les plages de lecture sont Noirs (0--10\,\%), Ombres (10--33\,\%),
Exposition (33--66\,\%), Hautes lumières (66--90\,\%) et Blancs (90--100\,\%).
Elles servent de repères visuels ; elles ne définissent pas des frontières
rigides d’action des curseurs.

\section{Les réglages}
\begin{longtable}{@{}L{3cm}L{2.8cm}L{7.8cm}@{}}
\toprule Réglage & Domaine & Utilisation\\\midrule\endhead
Exposition & $-5$ à $+5$ EV & Corriger globalement le niveau lumineux ; commencer par de petites variations.\\
Contraste & $-100$ à $+100$ & Renforcer ou adoucir la séparation des tons.\\
Ombres & $-100$ à $+100$ & Éclaircir ou assombrir les zones sombres sans cibler seulement l’extrémité noire.\\
Noirs & $-100$ à $+100$ & Ajuster les tons les plus sombres et l’assise de l’image.\\
Hautes lumières & $-100$ à $+100$ & Modérer ou renforcer les tons clairs.\\
Blancs & $-100$ à $+100$ & Ajuster les tons les plus proches du blanc.\\
\bottomrule
\end{longtable}
La valeur zéro est la référence de ces six réglages. La case \emph{Niveaux
de gris}, placée dans ce groupe, active ou désactive le rendu monochrome.

\section{Procédure conseillée}
Régler d’abord l’exposition, puis les ombres et hautes lumières, enfin les
noirs, blancs et le contraste. Revenir sur les étapes précédentes si nécessaire :
ces réglages interagissent visuellement. Une hausse des ombres rend aussi le
bruit présent plus visible. Un réglage de hautes lumières ne peut pas recréer
une information qui n’existe plus dans la source.

\paragraph{Exemple.}
Sur un contre-jour, essayer une légère baisse des hautes lumières et une
hausse modérée des ombres, sans déplacer immédiatement les deux curseurs
aux extrêmes. Maintenir \touche{Espace} dans la vue image pour contrôler
l’effet, puis vérifier les transitions au zoom avec \touche{Ctrl + 1}.
\gestesreglages

\chapter{Courbe de tonalité}
\section{Choisir un canal}
Choisir la luminance ou un canal rouge, vert ou bleu avec les boutons du module.
La diagonale représente la correspondance de référence. Une modification de
canal peut apporter une dominante colorée, même si la forme paraît proche
de celle de la courbe de luminance.

\section{Manipuler les points}
Double-cliquer dans le graphe pour ajouter une poignée, la faire glisser
pour modifier la courbe et cliquer avec le bouton droit sur une poignée
intérieure pour la supprimer. Les extrémités ne se manipulent pas comme les
points intérieurs. Éviter d’accumuler des points très proches : une courbe
plus simple est plus facile à contrôler.

Utiliser les commandes de réinitialisation de la courbe courante ou de
toutes les courbes selon l’intention. Le double-clic dans ce graphe
\emph{ajoute un point} : ce n’est pas le geste de remise à zéro des sliders.
Aucun raccourci spécifique d’ajout/suppression de points au clavier n’est déclaré.

\section{Le fond monochrome}
La silhouette grise est la luminance de l’aperçu courant ; elle reste monochrome
même si une courbe RGB est sélectionnée. C’est un repère de lecture après rendu,
pas l’histogramme exact du signal linéaire entrant dans les courbes.

\paragraph{Exemple.}
Pour donner un contraste doux, commencer par deux points intérieurs éloignés,
abaisser légèrement la zone sombre et relever légèrement la zone claire.
Comparer avec \touche{Ctrl + Maj + C} et surveiller les transitions plutôt
que chercher une forme géométrique parfaite.

\chapter{Couleur globale}
\section{Température et nuance}
La température couvre 2\,000 à 50\,000 K ; la nuance couvre $-150$ à $+150$.
La température déplace l’équilibre chaud/froid ; la nuance corrige l’axe
vert/magenta. Les bandes du curseur sont des indications, pas un nuancier
contractuel du résultat. La balance automatique se trouve dans le groupe RAW.

La température de référence dépend du développement de la photo ; elle n’est
pas toujours 5\,200 K ou 6\,500 K. La remise à zéro du contrôle de température
revient à sa température de référence, pas à zéro kelvin. La nuance se
réinitialise à zéro dans l’interface actuelle. Les unités ou valeurs affichées
par un autre logiciel ne garantissent pas une correspondance numérique exacte.

\section{Vibrance et saturation}
Les deux contrôles couvrent $-100$ à $+100$ avec une référence à zéro.
La saturation modifie la richesse colorée globale. La vibrance cherche une
action plus sélective suivant les couleurs déjà présentes ; ce n’est pas
simplement une saturation moins forte. Les effets peuvent devenir importants
sur les couleurs vives et modifier leur luminosité apparente.

\paragraph{Exemple.}
Commencer par une petite correction de température, puis régler la nuance.
Ajouter ensuite une vibrance modérée pour les couleurs ternes et vérifier
les carnations ainsi que les rouges et jaunes. N’ajuster la saturation
globale qu’en regardant l’ensemble de l’image et les zones déjà colorées.
Maintenir \touche{Espace} dans la vue pour comparer au rendu d’origine.

\attention{Les validations de couleur dépendent du rendu de départ, du profil
et du boîtier. Les valeurs extrêmes sont utiles pour tester, pas une promesse
de reproduction exacte de Lightroom.}
\gestesreglages

\chapter{Mélangeur de couleurs}
\section{Sélectionner une famille chromatique}
Les huit cercles sélectionnent les rouges, oranges, jaunes, verts, cyans,
bleus, violets ou magentas. Violet et magenta sont deux bandes distinctes.
Le nom coloré confirme la bande active. Changer de cercle change les
contrôles affichés, sans effacer les réglages des autres bandes.

\section{Teinte, saturation et luminance}
Chaque bande dispose de trois contrôles de $-100$ à $+100$, référence zéro.
La teinte déplace la famille vers ses voisines, la saturation modifie sa
richesse colorée et la luminance son niveau lumineux. Les familles se
recouvrent progressivement : une couleur située à une frontière peut
réagir à plusieurs bandes. Il ne s’agit pas d’une sélection d’objet.

\paragraph{Exemple.}
Pour travailler une fleur orange, sélectionner d’abord les oranges,
essayer une faible modification de luminance puis de saturation et examiner
les rouges et jaunes voisins. Si le rendu de départ est déjà trop rouge,
vérifier d’abord profil et balance des blancs plutôt que compenser systématiquement
ce défaut avec le mélangeur.

Comparer avant/après avec \touche{Ctrl + Maj + Y}. Aucun raccourci de sélection
des huit cercles n’est déclaré ; utiliser les cercles ou la navigation de focus.
\gestesreglages

\chapter{Color grading}
\section{Choisir une zone tonale}
Les trois boutons sélectionnent les ombres, les tons moyens et les hautes
lumières. Chaque zone conserve sa teinte et son intensité indépendamment des
deux autres. Contrairement au mélangeur, la sélection porte ici sur une zone
tonale, pas sur les objets déjà d’une certaine couleur.

\section{Utiliser la roue}
Cliquer dans la roue puis faire glisser la sélection. L’angle choisit une
teinte et l’éloignement du centre règle l’intensité. Pendant un glissement,
sortir du bord conserve la manipulation à l’intensité maximale ; il n’est
pas nécessaire de reprendre la poignée à chaque tour.
Le contrôle \emph{Intensité}, de 0 à 100, permet aussi une saisie numérique.
À zéro, la coloration de la zone est neutralisée même si une teinte reste mémorisée.

\paragraph{Exemple.}
Donner une faible teinte bleutée aux ombres et une faible teinte chaude aux
hautes lumières. Vérifier ensuite les tons moyens et les carnations, puis
comparer avec \touche{Espace} maintenu dans la vue. Éviter de pousser les
trois zones au maximum pour juger un rendu naturel.

\paragraph{Clavier.}
Le curseur d’intensité accepte les flèches, la saisie et le double-clic de
réinitialisation décrits pour les réglages. Aucun pilotage angulaire de la
roue par les flèches n’est actuellement déclaré.

\chapter{Effets}
\section{Choisir le type d’effet}
\begin{longtable}{@{}L{3.6cm}L{2.8cm}L{7.2cm}@{}}
\toprule Contrôle & Domaine & Intention\\\midrule\endhead
Texture & $-100$ à $+100$ & Adoucir ou renforcer les petits détails.\\
Clarté & $-100$ à $+100$ & Modifier la présence des structures et le contraste local.\\
Correction du voile & $-100$ à $+100$ & Modifier l’impression de voile et la séparation des tons.\\
Grain & 0 à 100 & Ajouter un effet de grain ; ce n’est pas du débruitage.\\
Flou gaussien & 0 à 20 & Adoucir globalement ; ce n’est pas une simulation complète de profondeur de champ.\\
\bottomrule
\end{longtable}
La référence de ces contrôles est zéro. Une valeur négative de texture ou
de clarté ne correspond pas à une réduction du bruit identique au module dédié.

\section{Contrôler les artefacts}
Examiner les contours, les transitions de ciel et les textures fines après
une forte correction. Un gain de clarté peut rendre le bruit plus visible ;
un excès de lissage peut donner une matière artificielle. Vérifier également
à la taille de livraison, car un détail agréable au fort zoom peut être excessif
en vue ajustée.

\paragraph{Exemple.}
Sur un paysage voilé, commencer par une petite correction du voile, puis
rééquilibrer éventuellement la lumière. Comparer avec \touche{Ctrl + Maj + Y}
et examiner les limites ciel/bâtiments au zoom \touche{Ctrl + 1}.
\gestesreglages

\chapter{Réduction du bruit}
\section{Séparer deux composantes}
\emph{Bruit de luminance} agit sur les fluctuations de clarté ;
\emph{Bruit de couleur} vise les variations chromatiques parasites.
Les deux contrôles vont de 0 à 100, référence zéro. Leur présence ne signifie
pas que le moteur utilise une IA équivalente à celle d’un autre logiciel.

\section{Méthode d’utilisation}
Corriger d’abord l’exposition utile et examiner une zone sombre relativement
uniforme, puis une zone texturée. Augmenter progressivement la réduction
de couleur si des taches colorées sont visibles. Ajuster ensuite la luminance
en vérifiant les cheveux, le tissu ou les détails fins. Revenir en vue ajustée
pour juger la sensation globale.

\paragraph{Exemple.}
Sur une photo à ISO élevés, ne pas conclure à partir du seul fond sombre.
Comparer également une texture et un contour net. Utiliser \touche{Ctrl + 1}
en gardant à l’esprit la résolution de l’aperçu, puis vérifier un export pleine
résolution pour une décision finale sur le détail.

\attention{Un débruitage plus fort peut supprimer de vrais détails.
Une référence Lightroom obtenue avec débruitage IA ne constitue pas une
comparaison à algorithmes identiques.}
\gestesreglages

\chapter{Netteté}
\section{Les trois contrôles}
Le \emph{Gain de netteté} va de 0 à 200 et règle la force du renforcement.
Le \emph{Rayon de netteté} va de 0,3 à 3,0 et règle son échelle ; sa référence
est 1,0. Le \emph{Masquage de netteté} va de 0 à 100 et module la sélectivité
du renforcement. Gain et masquage ont une référence à zéro.

\section{Procédure}
Évaluer d’abord le bruit, puis augmenter légèrement le gain de netteté.
Faire varier le rayon avec prudence et observer les contours contrastés.
Rechercher un gain de lisibilité sans liseré visible ni accentuation des
zones uniformes. Le renforcement ne récupère pas une mise au point réellement absente.

\paragraph{Exemple.}
Sur un portrait, contrôler le bord de l’œil et la texture de la peau, pas
uniquement les cheveux. Basculer entre \touche{Ctrl + 1} et \touche{Ctrl + 0}.
Maintenir \touche{Espace} pour comparer à la référence ; contrôler ensuite
un export à la taille finale.
\gestesreglages

\chapter{Corrections optiques}
\section{Activer les corrections disponibles}
Le groupe \emph{Corrections optiques} se trouve en fin de liste, avant Export.
Ses cases commandent la correction de l’objectif et la réduction de l’aberration
chromatique. Les capacités dépendent du RAW, des métadonnées et du profil
objectif disponible. Lire le diagnostic indiquant les corrections réellement
appliquées et l’éventuelle absence de profil LCP.

Le profil DCP du boîtier et le profil LCP de l’objectif ont des rôles différents.
Ne pas interpréter une case activée comme la preuve que toute correction
géométrique, de vignettage ou chromatique a effectivement pu être calculée.

\section{Comparer correctement}
Observer les lignes droites proches des bords, la luminosité des coins et
les franges sur les contours contrastés. Une correction de vignettage peut
éclaircir les bords et modifier l’impression lumineuse globale.
Pour comparer deux développements colorimétriques, garder les mêmes états
optiques ou les désactiver de manière cohérente des deux côtés.

\paragraph{Clavier.}
Utiliser \touche{Tab} pour atteindre une case puis \touche{Espace} pour la
basculer lorsqu’elle a le focus. Le même Espace n’active l’original temporaire
que dans la vue image. Il n’existe pas de raccourci global d’activation optique.

\part{Personnaliser et livrer}
\chapter{Préférences et disposition}
\section{Ouvrir les préférences}
Utiliser le menu \emph{Préférences} ou \touche{Ctrl + virgule}. Les trois onglets
sont \emph{Interface}, \emph{Apparence} et \emph{Avancé · RAW}.
Les préférences d’interface sont distinctes des retouches de chaque photo.

\section{Interface}
Afficher ou masquer Informations, Histogramme et Ruban ; choisir la largeur
du panneau Réglages (260--360 px) et la taille du texte des informations
(10--20 px). Les mêmes panneaux peuvent être masqués depuis le menu Affichage.

Les poignées entre les panneaux latéraux et la vue image permettent de
redimensionner Informations (210--300 px) et Réglages. Celle au-dessus du ruban
ajuste sa hauteur (100--260 px). Glisser à la souris, ou donner le focus à
la poignée et utiliser les flèches au pas de 10 px. Un double-clic restaure
la taille usuelle. Les dimensions sont mémorisées dans l’INI.

\section{Apparence}
Les palettes proposées sont Bordeaux, Blanc cassé, Noir mat, Bleu marine et
Gris actuel. Elles s’appliquent aux panneaux et aux préférences. Le fond de
la photo est indépendant : changer de palette ne remplace pas ce choix.
Les fonds bleu marine et rouge bordeaux utilisent des dégradés.

\attention{Pour juger finement la couleur d’une photo, un fond gris neutre
limite l’influence de l’entourage coloré. La palette de l’interface ne change
pas les pixels exportés et ne constitue pas un profil écran ICC.}

\section{Avancé · RAW}
L’onglet donne accès à la calibration de l’image ouverte, au manifeste et à
la réinitialisation profil/calibration. Ce dernier bouton est distinct de
\emph{Réinitialiser les réglages}, qui remet les retouches à leurs références
mais n’efface pas automatiquement tous les choix de développement RAW.
La réinitialisation générale est également accessible avec
\touche{Ctrl + Maj + R} et demande confirmation avant d’agir sur la photo ouverte.

\section{Persistance et état des réglages}
Les sections dépliées, dimensions, couleurs et certains choix d’interface
sont mémorisés. Le chemin du fichier INI apparaît dans les préférences.
Pour une édition manuelle, fermer IPR et conserver une copie du fichier.

Un nom de contrôle numérique en gras et un point sur sa section indiquent
une différence par rapport à sa valeur de réinitialisation. Ce repère ne
couvre pas encore toutes les courbes, cases et recettes ; ce n’est pas
un indicateur exhaustif de fichier enregistré ou non enregistré.
Les mémos explicatifs au survol prolongé restent prévus pour la phase bêta.

\chapter{Exporter une image}
\label{chap:export}
\section{Choisir un format}
Ouvrir la famille Export et choisir le format. Le contrôle Qualité JPEG
n’apparaît que pour JPEG ; il va de 1 à 100, référence 92.
\begin{longtable}{@{}L{2cm}L{4.6cm}L{7cm}@{}}
\toprule Format & Encodage actuel & Usage et limites\\\midrule\endhead
JPEG & sRGB, 8 bits & Copie courante ; compression avec perte, qualité réglable.\\
PNG & sRGB, 16 bits & Image rendue à haute profondeur ; pas de réglage Qualité JPEG.\\
TIFF & sRGB, 16 bits & Référence de comparaison ou échange de rendu à haute profondeur.\\
DNG & RGB linéaire, 16 bits, primaires sRGB & Contient le rendu développé ; ne préserve pas la mosaïque RAW d’origine.\\
\bottomrule
\end{longtable}
Ces propriétés décrivent l’exporteur actuel, pas toutes les capacités théoriques
de chaque format. Aucun sélecteur utilisateur de gamut de sortie ou de redimensionnement
n’est proposé dans ce panneau.

\section{Destination et confirmation}
\begin{enumerate}
\item Attendre la fin du développement et cliquer sur \emph{Choisir la destination…}.
\item Vérifier le nom, le dossier et le type du dialogue natif. Le format du
panneau est sélectionné initialement ; changer le type dans le dialogue change
aussi le format demandé.
\item Dans \emph{Vérifier l’export}, contrôler format, encodage, dimensions
pleine résolution et destination demandée.
\item Cliquer sur \emph{Exporter}, ou \emph{Annuler} pour revenir sans écrire.
\item Lire le message final pour connaître le fichier réellement publié ou l’erreur.
\end{enumerate}
L’extension est adaptée au format. Si le nom est pris, l’exporteur ajoute
un numéro disponible plutôt que d’écraser le fichier. Le chemin final peut
donc différer du nom demandé dans le récapitulatif.

\attention{Le récapitulatif ne promet pas un nom final réservé. Un autre
processus peut créer un fichier entre le choix et l’écriture ; la protection
contre les collisions reste assurée par l’exporteur.}

\section{Gestes et contrôles}
\touche{Tab} permet de parcourir les commandes du dialogue et \touche{Entrée}
de valider une saisie active. Utiliser explicitement le bouton Exporter pour
lancer l’écriture ; aucun \touche{Ctrl + E} global n’est déclaré dans IPR.
Ne pas confondre l’export avec la sauvegarde automatique des compagnons.

Pour un test reproductible, exporter une référence neutre et l’état retouché
avec le même profil, la même balance, les mêmes corrections optiques et le
même format. Noter les réglages : un TIFF IPR ne constitue pas à lui seul
un relevé complet et fiable de la recette appliquée.

\chapter{Vérifications et dépannage}
\section{Une commande paraît indisponible}
Vérifier qu’une photo est ouverte, attendre la fin du traitement et contrôler
qu’il s’agit d’un RAW pour une commande de développement appareil. Un profil
absent de la liste n’est pas nécessairement une erreur de la liste : sa
disponibilité dépend des fichiers installés et de la compatibilité.

\section{La photo semble différente d’un autre logiciel}
Comparer d’abord le même fichier, le même cadrage et le même état neutre.
Vérifier le profil appareil, la balance des blancs, la calibration, les corrections
optiques, le bruit et la netteté. Une même température affichée n’établit pas
à elle seule une égalité de développement. Le fond de visualisation et le
chemin de gestion couleur peuvent également influencer la comparaison.

\section{Des réglages réapparaissent à l’ouverture}
Rechercher le compagnon \chemin{.ipr} conservé avec la photo. Sa restauration
est volontaire. Utiliser la réinitialisation adaptée à l’intention, plutôt
que supprimer immédiatement des fichiers. Conserver une copie avant toute
intervention manuelle sur un compagnon ou une configuration.

\section{Une sauvegarde ou un export échoue}
Lire le message au bas de la fenêtre. Vérifier l’espace disque, les droits
du dossier, sa disponibilité et les éventuels verrouillages. Un message
d’échec de compagnon n’est pas un export réussi. Le message de réussite
d’export indique le chemin réel, y compris le numéro ajouté si nécessaire.

\section{Une fonction n’est pas décrite}
Les panneaux recadrage, masques et IA sont réservés dans cette édition.
Les badges complets du ruban, l’état exhaustif des groupes modifiés et les
états détaillés du moteur restent des évolutions, pas des commandes cachées.
La validation visuelle par l’utilisateur et les tests sur d’autres boîtiers
restent nécessaires malgré les tests automatisés de l’IHM.

\part{Familles réservées}
\chapter{Recadrage et effets géométriques}
% Famille réservée, comme dans la documentation technique. Aucun contenu tant que non commencée.

\chapter{Masques et retouche locale}
% Famille réservée, comme dans la documentation technique. Aucun contenu tant que non commencée.

\chapter{Intégration de l’intelligence artificielle}
% Famille réservée, comme dans la documentation technique. Aucun contenu tant que non commencée.

\appendix
\chapter{Aide-mémoire clavier et souris}
\label{chap:raccourcis}
Les gestes sont également expliqués dans leurs familles. Leur effet dépend
du focus : une flèche dans un curseur ajuste sa valeur ; ce n’est pas une
commande globale de changement de photo. Les raccourcis ci-dessous reprennent
les séquences déclarées dans la version décrite, notamment sur Windows.

\begin{longtable}{@{}L{4.4cm}L{9.2cm}@{}}
\toprule Touche ou geste & Action et contexte\\\midrule\endhead
\touche{Ctrl + virgule} & Ouvrir les préférences.\\
\touche{Ctrl + +} & Agrandir une photo prête. \touche{Ctrl + =} est également accepté.\\
\touche{Ctrl + $-$} & Réduire une photo prête.\\
\touche{Ctrl + 0} & Ajuster à la vue.\\
\touche{Ctrl + 1} & Afficher à 100\,\% de l’aperçu.\\
\touche{Ctrl + Maj + C} & Activer ou désactiver le comparatif partagé, si les deux aperçus sont disponibles.\\
\touche{Espace} maintenu & Original temporaire lorsque la vue image a le focus ; relâcher restaure la vue.\\
\touche{Ctrl + molette} & Zoom centré autour du pointeur dans la vue image.\\
Molette sans Ctrl & Déplacer l’image si elle dépasse la vue.\\
\touche{Maj + molette} & Déplacement horizontal à partir de la molette verticale, sur une image agrandie.\\
Clic dans la photo & Basculer entre ajustement et agrandissement, sans glissement.\\
Clic gauche maintenu, déplacement & Parcourir une image agrandie ; ou déplacer le séparateur du comparatif.\\
\bottomrule
\end{longtable}

\clearpage
\section*{Réglages et manipulation des contrôles}
\begin{longtable}{@{}L{4.4cm}L{9.2cm}@{}}
\toprule Touche ou geste & Action et contexte\\\midrule\endhead
\touche{Tab}, \touche{Maj + Tab} & Parcourir les contrôles qui acceptent le focus.\\
\touche{Espace} sur un bouton ou une case & Activer la commande focalisée ; ne pas confondre avec le maintien dans la vue image.\\
Flèches sur un slider & Ajuster au pas déclaré : par exemple 0,01 EV pour l’exposition.\\
\touche{Entrée} dans un champ & Confirmer la saisie numérique.\\
Double-clic sur nom ou slider & Restaurer la référence du paramètre ; cette référence n’est pas toujours zéro.\\
Flèches sur une poignée de panneau & Redimensionner au pas de 10 px dans sa direction.\\
Double-clic sur une poignée de panneau & Restaurer sa dimension usuelle.\\
Double-clic dans les courbes & Ajouter un point de contrôle.\\
Clic droit sur un point intérieur & Supprimer ce point de courbe.\\
Glisser dans une roue de grading & Modifier teinte et intensité de la zone active.\\
\bottomrule
\end{longtable}

\attention{Ne pas supposer les raccourcis de Lightroom ou d’un autre éditeur :
Ctrl + O, Ctrl + S, Ctrl + E, Ctrl + Z et les touches numériques de notation
ne sont pas déclarés comme commandes globales de l’application décrite.
Les raccourcis propres aux dialogues natifs dépendent du système.}

\chapter{Glossaire pratique}
\begin{description}
\item[RAW] Données de prise de vue à développer ; leur rendu dépend notamment
du décodage, de la balance des blancs et des profils.
\item[DCP] Profil couleur appareil utilisé dans le développement RAW.
\item[LCP] Profil objectif permettant certaines corrections optiques.
\item[ICC] Description colorimétrique d’un espace ou d’un périphérique ;
distincte d’une palette de couleurs de l’interface.
\item[EV ou IL] Unité de variation d’exposition ; une augmentation de 1 EV
correspond à un gain lumineux linéaire de facteur deux avant le rendu de sortie.
\item[Température] Réglage de balance des blancs exprimé en kelvins, à
distinguer d’une mesure universelle de tous les pixels de l’image.
\item[Nuance] Correction suivant l’axe vert/magenta.
\item[Teinte] Famille chromatique, par exemple rouge ou orange ; ne pas
confondre avec le nom du curseur Nuance.
\item[Luminance] Grandeur relative décrivant le niveau lumineux ; sa définition
précise dépend de l’espace et de l’étape de calcul.
\item[Gamut] Ensemble des couleurs représentables dans un espace donné.
\item[sRGB] Espace de sortie courant des exports rendus décrits dans ce livre.
\item[Aperçu] Image calculée pour la visualisation, éventuellement moins définie
que la source ; ce n’est pas la preuve de résolution de l’export.
\item[Compagnon ou sidecar] Fichier associé à la source, ici \chemin{.ipr},
conservant des réglages sans remplacer les pixels du fichier source.
\item[INI] Fichier des préférences d’interface, indépendant des compagnons photo.
\item[Neutre] Référence de réglages déclarée pour une comparaison ; elle doit
toujours préciser le développement de départ, le profil et la balance des blancs.
\end{description}

\end{document}
